\section{Optimality of the Converged Solution}

In this section, we connect the optimality of the solution obtained from \cref{alg:sqrt_qr_covariance_steering} (in the $\mathbb{L}_+^n$ coordinates) to that obtained from solving it defining the covariance matrix as variables, as in \cite{liuOptimalCovarianceSteering2025} (in the $\mathbb{S}_+^n$ coordinates).

 Consider two constrained optimization problems on manifolds $\mathcal{M}_A$ and $\mathcal{M}_B$ with variables $x \in \mathcal{M}_A$ and $y \in \mathcal{M}_B$:%
\begin{equation*}
\text{(A)} \\ \minimize_{x \in \mathcal{X}} \ f(x) \qquad
\text{(B)} \\ \minimize_{y \in \mathcal{Y}} \ f(T(y))
\end{equation*}
where $T: \mathcal{M}_B \rightarrow \mathcal{M}_A$ is a diffeomorphism such that $T(\mathcal{M}_B) = \mathcal{M}_A$, with equality constraints $g_i, i \in \mathcal{I}$ and inequality constraints $h_j, j \in \mathcal{J}$:
\begin{align}
	\mathcal{X} &:= \{ x \in \mathcal{M}_A \ |\  g_i(x) = 0, i \in \mathcal{I}; \ h_j(x) \leq 0 \in \mathcal{J} \} \nonumber \\ 
	\mathcal{Y} &:= \{ y \in \mathcal{M}_B \ |\  g_i(T(y)) = 0, i \in \mathcal{I}; \ h_j(T(y)) \leq 0\in \mathcal{J} \} \nonumber 
\end{align}
From the definition, $\mathcal{X} = T(\mathcal{Y})$.

\begin{proposition}[Invariance of Local Optimality via a Diffeomorphism]
\label{prop:global_optimality_diffeomorphism}
Let $y^*$ be a local minimum of (B). Then $x^* = T(y^*)$ is a local minimum of (A).
\end{proposition}

\begin{proof}
Since $y^*$ is a local minimum of (B), there exists a neighborhood $\mathcal{N}$ around $y^*$ such that
\begin{equation}
\label{eq:local_minimum}
f(T(y)) \geq f(T(y^*)) \quad \forall y \in \mathcal{S}_B
\end{equation}
where $\mathcal{S}_B := N \cap \mathcal{Y}$ \cite[Chapter 12]{nocedalNumericalOptimization2006}.

Since $T$ is a diffeomorphism, it is in particular a homeomorphism. Therefore, $T$ maps the open neighborhood $\mathcal{N}$ to an open set $T(\mathcal{N})$ which contains $T(y^*) = x^*$ \cite[Chapter 12]{munkresTopology2014}.
Defining $\mathcal{S}_A := T(N) \cap \mathcal{X}$,
\begin{equation}
\mathcal{S}_A = T(\mathcal{N}) \cap T(\mathcal{Y}) 
= T(\mathcal{N} \cap \mathcal{Y}) 
= T(\mathcal{S}_B)
\end{equation}
where the second equality holds since $T$ is in particular a bijection \cite[Theorem 12.4]{hammackBookProof2018}.
Therefore, for any $x \in \mathcal{S}_A$, there exists a unique pre-image $y = T^{-1}(x) \in \mathcal{S}_B$. From \cref{eq:local_minimum},
\begin{equation}
f(T(y)) \geq f(T(y^*)),\ \forall y \in \mathcal{S}_B \nonumber
\Rightarrow f(x) \geq f(x^*),\ \forall x \in \mathcal{S}_A. \label{eq:local_optimality_x}
\end{equation}
Hence, \(x^*\) is a local minimum of (A).
\end{proof}

\begin{remark}
	Another way to show \cref{prop:global_optimality_diffeomorphism} is to show that the KKT conditions of (A) and (B) are equivalent. Readers are referred to Proposition 1 of \cite{fatkhullinGlobalSolutionsNonConvex2025}. The assumptions used in the proposition, 1. weak convexity of the objective and constraints and 2. differentiability of the coordinate transformation, are satisfied by the problem statement in \cref{eq:nonconvex_problem_statement}.
\end{remark}

Next, we apply this result to our specific problem. The following proposition is used in the proof.
\begin{proposition}[\cite{linRiemannianGeometrySymmetric2019}]
	\label{prop:cholesky_diffeomorphism}
	The Cholesky map \(\mathcal{L} : \mathbb{S}^n_+ \rightarrow \mathbb{L}_+^n\) is a diffeomorphism.
\end{proposition}

\begin{theorem} \label{th:local_optimality}
	\cref{alg:sqrt_qr_covariance_steering} converges to a locally optimal solution of the cc-CS problem. The local optimality is invariant to whether the problem is represented with the full covariance matrix or its square root. 
\end{theorem}

\begin{proof}
	Let (A) from \cref{prop:global_optimality_diffeomorphism} be the determinisic formulation method in \cite{liuOptimalCovarianceSteering2025}. In this formulation, variable $x \in \mathcal{M}_A = \mathbb{S}_+^n \times \cdots \times \mathbb{S}_+^n \times \R^{n\times m} \times \cdots \times \R^{n \times m}$ where the first $N+1$ spaces are the PD matrices $P_k$ and the next $N$ spaces are the feedback gains $K_k$. 
	Let (B) from \cref{prop:global_optimality_diffeomorphism} be \cref{eq:nonconvex_problem_statement} before the change of variables $L_k = K_k S_k$. In this formulation, variable $y \in \mathcal{M}_B = \mathbb{L}_+^n \times \cdots \times \mathbb{L}_+^n \times \R^{n\times m} \times \cdots \times \R^{n \times m}$, where the first $N+1$ spaces are $S_k$ and the next $N$ spaces are the feedback gains $K_k$. (We can ignore the mean variables $\bm{v}, \bm{\mu}$ as they appear exactly in the same way in either problem.)
	(B) is solved via \cref{alg:sqrt_qr_covariance_steering}, and the diffeomorphism $T: \mathcal{M}_B \rightarrow \mathcal{M}_A$ maps each of the matrices $\mathbb{S}_+^n$ to $\mathbb{L}_+^n$ via the inverse of the Cholesky map \(\mathcal{L} : \mathbb{S}^n_+ \rightarrow \mathbb{L}_+^n\) and does not modify the feedback gains $K_k$. The inverse of a diffeomorphism is also a diffeomorphism.
	Then, the proof follows from the local optimality of the \texttt{SCvx*} algorithm \cite[Theorem 3]{oguriSuccessiveConvexificationFeasibility2023a} and \cref{prop:cholesky_diffeomorphism}. 
\end{proof}


\begin{theorem} \label{th:global_optimality}
	When the CCs are not present and the objective takes the quadratic form of \cref{eq:objective_function}, the solution obtained by \cref{alg:sqrt_qr_covariance_steering} is globally optimal.
\end{theorem}
\begin{proof}
	The unconstrained CS problem with quadratic cost is shown to be a convex optimization problem \cite{liuOptimalCovarianceSteering2025}. Since a local optimum of a convex problem is also a global optimum, from \cref{th:local_optimality}, the locally optimal solution obtained from \cref{alg:sqrt_qr_covariance_steering} is also globally optimal.
\end{proof}